% Options for packages loaded elsewhere
\PassOptionsToPackage{unicode}{hyperref}
\PassOptionsToPackage{hyphens}{url}
\documentclass[
]{article}
\usepackage{xcolor}
\usepackage[margin=1in]{geometry}
\usepackage{amsmath,amssymb}
\setcounter{secnumdepth}{-\maxdimen} % remove section numbering
\usepackage{iftex}
\ifPDFTeX
  \usepackage[T1]{fontenc}
  \usepackage[utf8]{inputenc}
  \usepackage{textcomp} % provide euro and other symbols
\else % if luatex or xetex
  \usepackage{unicode-math} % this also loads fontspec
  \defaultfontfeatures{Scale=MatchLowercase}
  \defaultfontfeatures[\rmfamily]{Ligatures=TeX,Scale=1}
\fi
\usepackage{lmodern}
\ifPDFTeX\else
  % xetex/luatex font selection
\fi
% Use upquote if available, for straight quotes in verbatim environments
\IfFileExists{upquote.sty}{\usepackage{upquote}}{}
\IfFileExists{microtype.sty}{% use microtype if available
  \usepackage[]{microtype}
  \UseMicrotypeSet[protrusion]{basicmath} % disable protrusion for tt fonts
}{}
\makeatletter
\@ifundefined{KOMAClassName}{% if non-KOMA class
  \IfFileExists{parskip.sty}{%
    \usepackage{parskip}
  }{% else
    \setlength{\parindent}{0pt}
    \setlength{\parskip}{6pt plus 2pt minus 1pt}}
}{% if KOMA class
  \KOMAoptions{parskip=half}}
\makeatother
\usepackage{color}
\usepackage{fancyvrb}
\newcommand{\VerbBar}{|}
\newcommand{\VERB}{\Verb[commandchars=\\\{\}]}
\DefineVerbatimEnvironment{Highlighting}{Verbatim}{commandchars=\\\{\}}
% Add ',fontsize=\small' for more characters per line
\usepackage{framed}
\definecolor{shadecolor}{RGB}{248,248,248}
\newenvironment{Shaded}{\begin{snugshade}}{\end{snugshade}}
\newcommand{\AlertTok}[1]{\textcolor[rgb]{0.94,0.16,0.16}{#1}}
\newcommand{\AnnotationTok}[1]{\textcolor[rgb]{0.56,0.35,0.01}{\textbf{\textit{#1}}}}
\newcommand{\AttributeTok}[1]{\textcolor[rgb]{0.13,0.29,0.53}{#1}}
\newcommand{\BaseNTok}[1]{\textcolor[rgb]{0.00,0.00,0.81}{#1}}
\newcommand{\BuiltInTok}[1]{#1}
\newcommand{\CharTok}[1]{\textcolor[rgb]{0.31,0.60,0.02}{#1}}
\newcommand{\CommentTok}[1]{\textcolor[rgb]{0.56,0.35,0.01}{\textit{#1}}}
\newcommand{\CommentVarTok}[1]{\textcolor[rgb]{0.56,0.35,0.01}{\textbf{\textit{#1}}}}
\newcommand{\ConstantTok}[1]{\textcolor[rgb]{0.56,0.35,0.01}{#1}}
\newcommand{\ControlFlowTok}[1]{\textcolor[rgb]{0.13,0.29,0.53}{\textbf{#1}}}
\newcommand{\DataTypeTok}[1]{\textcolor[rgb]{0.13,0.29,0.53}{#1}}
\newcommand{\DecValTok}[1]{\textcolor[rgb]{0.00,0.00,0.81}{#1}}
\newcommand{\DocumentationTok}[1]{\textcolor[rgb]{0.56,0.35,0.01}{\textbf{\textit{#1}}}}
\newcommand{\ErrorTok}[1]{\textcolor[rgb]{0.64,0.00,0.00}{\textbf{#1}}}
\newcommand{\ExtensionTok}[1]{#1}
\newcommand{\FloatTok}[1]{\textcolor[rgb]{0.00,0.00,0.81}{#1}}
\newcommand{\FunctionTok}[1]{\textcolor[rgb]{0.13,0.29,0.53}{\textbf{#1}}}
\newcommand{\ImportTok}[1]{#1}
\newcommand{\InformationTok}[1]{\textcolor[rgb]{0.56,0.35,0.01}{\textbf{\textit{#1}}}}
\newcommand{\KeywordTok}[1]{\textcolor[rgb]{0.13,0.29,0.53}{\textbf{#1}}}
\newcommand{\NormalTok}[1]{#1}
\newcommand{\OperatorTok}[1]{\textcolor[rgb]{0.81,0.36,0.00}{\textbf{#1}}}
\newcommand{\OtherTok}[1]{\textcolor[rgb]{0.56,0.35,0.01}{#1}}
\newcommand{\PreprocessorTok}[1]{\textcolor[rgb]{0.56,0.35,0.01}{\textit{#1}}}
\newcommand{\RegionMarkerTok}[1]{#1}
\newcommand{\SpecialCharTok}[1]{\textcolor[rgb]{0.81,0.36,0.00}{\textbf{#1}}}
\newcommand{\SpecialStringTok}[1]{\textcolor[rgb]{0.31,0.60,0.02}{#1}}
\newcommand{\StringTok}[1]{\textcolor[rgb]{0.31,0.60,0.02}{#1}}
\newcommand{\VariableTok}[1]{\textcolor[rgb]{0.00,0.00,0.00}{#1}}
\newcommand{\VerbatimStringTok}[1]{\textcolor[rgb]{0.31,0.60,0.02}{#1}}
\newcommand{\WarningTok}[1]{\textcolor[rgb]{0.56,0.35,0.01}{\textbf{\textit{#1}}}}
\usepackage{graphicx}
\makeatletter
\newsavebox\pandoc@box
\newcommand*\pandocbounded[1]{% scales image to fit in text height/width
  \sbox\pandoc@box{#1}%
  \Gscale@div\@tempa{\textheight}{\dimexpr\ht\pandoc@box+\dp\pandoc@box\relax}%
  \Gscale@div\@tempb{\linewidth}{\wd\pandoc@box}%
  \ifdim\@tempb\p@<\@tempa\p@\let\@tempa\@tempb\fi% select the smaller of both
  \ifdim\@tempa\p@<\p@\scalebox{\@tempa}{\usebox\pandoc@box}%
  \else\usebox{\pandoc@box}%
  \fi%
}
% Set default figure placement to htbp
\def\fps@figure{htbp}
\makeatother
\setlength{\emergencystretch}{3em} % prevent overfull lines
\providecommand{\tightlist}{%
  \setlength{\itemsep}{0pt}\setlength{\parskip}{0pt}}
\usepackage{bookmark}
\IfFileExists{xurl.sty}{\usepackage{xurl}}{} % add URL line breaks if available
\urlstyle{same}
\hypersetup{
  pdftitle={Rukun's work},
  hidelinks,
  pdfcreator={LaTeX via pandoc}}

\title{Rukun's work}
\author{}
\date{\vspace{-2.5em}2026-08-18}

\begin{document}
\maketitle

\subsection{1.Study population
characteristics}\label{study-population-characteristics}

Before modelling the association between greenspace and asthma, we first
describe the participating populations.

\paragraph{What does the distribution of variables look
like?}\label{what-does-the-distribution-of-variables-look-like}

--\textgreater{} Heatmap We can first examine lsoa level summary
statistics : 1. mean 2. Std deviation 3. Max 4. Min

Stats can be split/combined across study option to return plot/table

\begin{Shaded}
\begin{Highlighting}[]
\FunctionTok{ds.geoheatmapPlot}\NormalTok{(}\AttributeTok{x=}\StringTok{"D$has\_asthma"}\NormalTok{, }\AttributeTok{y =} \StringTok{"D$lsoa11cd"}\NormalTok{, }\AttributeTok{datasources =}\NormalTok{ conns)}
\FunctionTok{ds.geoheatmapPlot}\NormalTok{(}\AttributeTok{x=}\StringTok{"D$has\_asthma"}\NormalTok{, }\AttributeTok{y =} \StringTok{"D$lsoa11cd"}\NormalTok{, }\AttributeTok{datasources =}\NormalTok{ conns,}
\AttributeTok{type=}\StringTok{\textquotesingle{}split\textquotesingle{}}\NormalTok{)}
\FunctionTok{ds.geoheatmapPlot}\NormalTok{(}\AttributeTok{x=}\StringTok{"D$has\_asthma"}\NormalTok{, }\AttributeTok{y =} \StringTok{"D$lsoa11cd"}\NormalTok{, }\AttributeTok{datasources =}\NormalTok{ conns,}
\AttributeTok{type=}\StringTok{\textquotesingle{}split\textquotesingle{}}\NormalTok{, }\AttributeTok{names\_region =} \FunctionTok{c}\NormalTok{(}\StringTok{"Liverpool"}\NormalTok{, }\StringTok{"Sefton"}\NormalTok{))}
\FunctionTok{ds.geoheatmapPlot}\NormalTok{(}\AttributeTok{x=}\StringTok{"D$distance\_local\_greenspace"}\NormalTok{, }\AttributeTok{y =} \StringTok{"D$lsoa11cd"}\NormalTok{,}
\AttributeTok{datasources =}\NormalTok{ conns)}
\FunctionTok{ds.geoheatmapPlot}\NormalTok{(}\AttributeTok{x=}\StringTok{"D$has\_asthma"}\NormalTok{, }\AttributeTok{y =} \StringTok{"D$lsoa11cd"}\NormalTok{, }\AttributeTok{datasources =}\NormalTok{ conns,}
\AttributeTok{type =} \StringTok{\textquotesingle{}split\textquotesingle{}}\NormalTok{, }\AttributeTok{names\_region =} \FunctionTok{c}\NormalTok{(}\StringTok{"Liverpool"}\NormalTok{, }\StringTok{"Sefton"}\NormalTok{))}
\end{Highlighting}
\end{Shaded}

---\textgreater{} Histogram : may not be useful for lsoa level analysis

\paragraph{Is greenspace distance approximately normally
distributed?}\label{is-greenspace-distance-approximately-normally-distributed}

This is to be developed

\begin{Shaded}
\begin{Highlighting}[]
\FunctionTok{ds.qqnorm}\NormalTok{()}
\end{Highlighting}
\end{Shaded}

\paragraph{Does asthma prevalence differ by
sex/age/ethnicity/IMD?}\label{does-asthma-prevalence-differ-by-sexageethnicityimd}

For two categorical variables, we can use a chi-squared test.

\begin{Shaded}
\begin{Highlighting}[]
\FunctionTok{ds.chisqr}\NormalTok{(}\AttributeTok{rvar=}\StringTok{"D$sex"}\NormalTok{, }\AttributeTok{cvar=}\StringTok{"D$has\_asthma"}\NormalTok{, }\AttributeTok{stvar=}\StringTok{"D$lsoa11cd"}\NormalTok{, }
          \AttributeTok{exclude=}\ConstantTok{NULL}\NormalTok{, }\AttributeTok{useNA =}\StringTok{"no"}\NormalTok{, }\AttributeTok{suppress.chisq.warnings=}\ConstantTok{FALSE}\NormalTok{,}
          \AttributeTok{datasources=}\ConstantTok{NULL}\NormalTok{, }\AttributeTok{names\_region =} \FunctionTok{c}\NormalTok{(}\StringTok{"Liverpool"}\NormalTok{,  }\StringTok{"Sefton"}\NormalTok{),}
          \AttributeTok{force.nfilter=}\ConstantTok{NULL}\NormalTok{, }\AttributeTok{split =} \ConstantTok{FALSE}\NormalTok{, }\AttributeTok{draw.plot =} \ConstantTok{TRUE}\NormalTok{)}

\FunctionTok{ds.chisqr}\NormalTok{(}\AttributeTok{rvar=}\StringTok{"D$sex"}\NormalTok{, }\AttributeTok{cvar=}\StringTok{"D$has\_asthma"}\NormalTok{, }\AttributeTok{stvar=}\StringTok{"D$lsoa11cd"}\NormalTok{, }
          \AttributeTok{exclude=}\ConstantTok{NULL}\NormalTok{, }\AttributeTok{useNA =}\StringTok{"no"}\NormalTok{, }\AttributeTok{suppress.chisq.warnings=}\ConstantTok{FALSE}\NormalTok{,}
          \AttributeTok{datasources=}\ConstantTok{NULL}\NormalTok{, }\AttributeTok{names\_region =} \StringTok{"Liverpool"}\NormalTok{,}
          \AttributeTok{force.nfilter=}\ConstantTok{NULL}\NormalTok{, }\AttributeTok{split =} \ConstantTok{TRUE}\NormalTok{, }\AttributeTok{draw.plot =} \ConstantTok{TRUE}\NormalTok{)}


\FunctionTok{ds.chisqr}\NormalTok{(}\AttributeTok{rvar=}\StringTok{"D$sex"}\NormalTok{, }\AttributeTok{cvar=}\StringTok{"D$has\_asthma"}\NormalTok{, }\AttributeTok{stvar=}\StringTok{"D$lsoa11cd"}\NormalTok{,}
          \AttributeTok{exclude=}\ConstantTok{NULL}\NormalTok{, }\AttributeTok{useNA =}\StringTok{"no"}\NormalTok{, }\AttributeTok{suppress.chisq.warnings=}\ConstantTok{FALSE}\NormalTok{,}
          \AttributeTok{datasources=}\ConstantTok{NULL}\NormalTok{, }\AttributeTok{names\_region =} \ConstantTok{NULL}\NormalTok{,}
          \AttributeTok{force.nfilter=}\ConstantTok{NULL}\NormalTok{, }\AttributeTok{split =} \ConstantTok{TRUE}\NormalTok{, }\AttributeTok{draw.plot =} \ConstantTok{FALSE}\NormalTok{)}

\FunctionTok{ds.chisqr}\NormalTok{(}\AttributeTok{rvar=}\StringTok{"D$sex"}\NormalTok{, }\AttributeTok{cvar=}\StringTok{"D$has\_asthma"}\NormalTok{, }\AttributeTok{stvar=}\StringTok{"D$lsoa11cd"}\NormalTok{,}
          \AttributeTok{exclude=}\ConstantTok{NULL}\NormalTok{, }\AttributeTok{useNA =}\StringTok{"no"}\NormalTok{, }\AttributeTok{suppress.chisq.warnings=}\ConstantTok{FALSE}\NormalTok{,}
          \AttributeTok{datasources=}\ConstantTok{NULL}\NormalTok{, }\AttributeTok{names\_region =} \ConstantTok{NULL}\NormalTok{,}
          \AttributeTok{force.nfilter=}\ConstantTok{NULL}\NormalTok{, }\AttributeTok{split =} \ConstantTok{TRUE}\NormalTok{, }\AttributeTok{draw.plot =} \ConstantTok{TRUE}\NormalTok{)}

\FunctionTok{ds.chisqr}\NormalTok{(}\AttributeTok{rvar=}\StringTok{"D$sex"}\NormalTok{, }\AttributeTok{cvar=}\StringTok{"D$has\_asthma"}\NormalTok{, }\AttributeTok{stvar=}\StringTok{"D$lsoa11cd"}\NormalTok{,}
          \AttributeTok{exclude=}\ConstantTok{NULL}\NormalTok{, }\AttributeTok{useNA =}\StringTok{"no"}\NormalTok{, }\AttributeTok{suppress.chisq.warnings=}\ConstantTok{FALSE}\NormalTok{,}
          \AttributeTok{datasources=}\ConstantTok{NULL}\NormalTok{, }\AttributeTok{names\_region =} \StringTok{"Liverpool"}\NormalTok{,}
          \AttributeTok{force.nfilter=}\ConstantTok{NULL}\NormalTok{, }\AttributeTok{split =} \ConstantTok{TRUE}\NormalTok{, }\AttributeTok{draw.plot =} \ConstantTok{FALSE}\NormalTok{)}
\end{Highlighting}
\end{Shaded}

age/continous variable can't be tested here (develop t-test?)

\subsection{2. Effect of distance to
greenspace}\label{effect-of-distance-to-greenspace}

\paragraph{Is greenspace distance associated with
asthma?}\label{is-greenspace-distance-associated-with-asthma}

This is a function from dsbase.

\begin{Shaded}
\begin{Highlighting}[]
\FunctionTok{ds.glm}\NormalTok{(}\AttributeTok{formula=}\StringTok{"has\_asthma \textasciitilde{} greenspace\_distance"}\NormalTok{,}\AttributeTok{family=}\StringTok{"binomial"}\NormalTok{,}\AttributeTok{datasources=}\NormalTok{connections}
\NormalTok{)}
\end{Highlighting}
\end{Shaded}

\paragraph{Does the association remain after accounting for individual
characteristics?}\label{does-the-association-remain-after-accounting-for-individual-characteristics}

This is a function from dsbase.

\begin{Shaded}
\begin{Highlighting}[]
\FunctionTok{ds.glm}\NormalTok{(}\AttributeTok{formula=}\FunctionTok{paste}\NormalTok{(}\StringTok{"has\_asthma \textasciitilde{} greenspace\_distance +"}\NormalTok{,}\StringTok{"age + sex + ethnicity + imd\_decile"}
\NormalTok{  ),}\AttributeTok{family=}\StringTok{"binomial"}\NormalTok{,}\AttributeTok{datasources=}\NormalTok{connections}
\NormalTok{)}
\end{Highlighting}
\end{Shaded}

\subsubsection{3. Accounting for geographical
clustering}\label{accounting-for-geographical-clustering}

\paragraph{Question 3.1 --- Is greenspace associated with asthma after
accounting for LSOA
clustering?}\label{question-3.1-is-greenspace-associated-with-asthma-after-accounting-for-lsoa-clustering}

This is a function from dsbase which is based on glmer: a function from
package lme4.

\begin{Shaded}
\begin{Highlighting}[]
\FunctionTok{ds.glmerSLMA}\NormalTok{(}\AttributeTok{formula=}\StringTok{"has\_asthma \textasciitilde{} greenspace\_distance + (1|lsoa11cd)"}\NormalTok{,}\AttributeTok{family=}\StringTok{"binomial"}\NormalTok{,}\AttributeTok{datasources=}\NormalTok{connections}
\NormalTok{)}
\end{Highlighting}
\end{Shaded}


\end{document}
